
\section{从\trLocallyRingedSpace{}到概形的态射}

为了证明\trMorphism{}确会如我们所望的那样行事，我们将动用\trLocallyRingedSpace{}这一观念。虽然它在几何学的其他场合自有用处，但往后我们不会再用到它。\trLocallyRingedSpace（连同其间的映射）这一观念，其灵感来自我们对\trManifold{}的已有认识（见\cref{ex:stalk-map-preserves-vanishing}）。设 $\pi\colon X\to Y$ 是\trManifold{}之间的一个\trMorphism，$\pi(p)=q$，而 $f$ 是 $Y$ 上一个在 $q$ 处为零的函数，那么拉回到 $X$ 上的函数 $\pi^{\sharp}(f)$ 就应当在 $p$ 处为零。换个说法：在 $q$ 处为零的函数\trGerm（于 $q\in Y$ 处取的），拉回后应当成为在 $p$ 处为零的函数\trGerm（于 $p\in X$ 处取的）。

\textbf{7.3.1.}\ \emph{\trhDefinition{}。}\quad 回想（定义~4.3.7）：\trLocallyRingedSpace{}是这样的\trRingedSpace{} $(X,\OO_X)$，其诸\trStalk{} $\OO_{X,p}$ 全都是\trLocalRing。\textbf{\trLocallyRingedSpace{}之间的\trMorphism} $\pi\colon X\to Y$，是这样一个\trRingedSpace{}的\trMorphism：它所诱导的诸\trStalk{}之间的映射 $\pi^{\sharp}\colon\OO_{Y,q}\to\OO_{X,p}$（练习~7.2.E）把 $\OO_{Y,q}$ 的\trMaximalIdeal{} $\mathfrak{m}_{Y,q}$ 送入 $\OO_{X,p}$ 的\trMaximalIdeal{} $\mathfrak{m}_{X,p}$（即一个「\trLocalRing{}的\trRingMorphism」（morphism of local rings））。（特别提请注意：\trLocallyRingedSpace{}构成一个\trCategory。）这话说的是一件相当具体而直观的事情：「若 $p\mapsto q$，而 $g$ 是一个在 $q$ 处为零的函数，那么它拉回后便是一个在 $p$ 处为零的函数。」你还会想要：「若 $p\mapsto q$，而 $g$ 是一个在 $q$ 处\emph{不}为零的函数，那么它拉回后便是一个在 $p$ 处\emph{不}为零的函数。」而这正是下一道练习的一个推论。

\textbf{7.3.A.（练习：\trLocallyRingedSpace{}上函数的值可拉回）}\quad 设 $\pi\colon X\to Y$ 是\trLocallyRingedSpace{}之间的一个\trMorphism，且 $\pi(p)=q$。证明：$\pi$ 诱导诸\trResidueField{}之间的一个包含映射 $\kappa(q)\hookrightarrow\kappa(p)$。这就抓住了如下事实：拉回把「函数为零（相应地，非零）的\trLocus」送到「拉回后的函数为零（相应地，非零）的\trLocus」。特别地，由\cref{ex:vanishing-locus-closed}(b)，可逆函数的拉回仍是可逆函数。

小结一下：\trLocallyRingedSpace{}这个观念，我们只在定义\trScheme{}的\trMorphism{}并证明\trMorphism{}具有种种合理性质时才会用到。关于\trLocallyRingedSpace，你需要记住的要点有二：(i) 函数在点处有值；(ii) 给定\trLocallyRingedSpace{}之间的一个映射，函数的值会「拉回」，而且特别地，函数的零点也会「拉回」。

\textbf{7.3.B.（练习）}\quad 证明：\trLocallyRingedSpace{}的\trMorphism{}可粘（对照练习~7.2.A）。（提示：你对练习~7.2.A 的解答或许原封不动就能用。）

\textbf{7.3.C.（容易的重要练习）}\quad 回想（\S4.3.7）：$\Spec A$ 是一个\trLocallyRingedSpace。证明：由\trRingMorphism{} $\pi^{\sharp}\colon B\to A$ 所定义的\trRingedSpace{}的\trMorphism{} $\pi\colon\Spec A\to\Spec B$（练习~7.2.F）是\trLocallyRingedSpace{}之间的一个\trMorphism。

\textbf{7.3.2.}\ \emph{\textbf{关键命题。}}\ ——\ \emph{若 $\pi\colon\Spec A\to\Spec B$ 是\trLocallyRingedSpace{}之间的一个\trMorphism，那么它正是由映射 $\pi^{\sharp}\colon B=\Gamma(\Spec B,\OO_{\Spec B})\to\Gamma(\Spec A,\OO_{\Spec A})=A$ 依练习~7.3.C 之例所诱导的那个\trLocallyRingedSpace{}之间的\trMorphism。}

（旁注：\cref{ex:scheme-isos-ring-isos} 是关键命题~7.3.2 的一个特例。你应当回头看看自己对\cref{ex:scheme-isos-ring-isos}的解答，找一找自己在何处暗中用到了关于\trLocallyRingedSpace{}的想法。）

\emph{证明。}\ 设 $\pi\colon\Spec A\to\Spec B$ 是\trLocallyRingedSpace{}之间的一个\trMorphism。我们要证明：它被其在\trGlobalSections{}上的映射 $\pi^{\sharp}\colon B\to A$ 完全确定。首先需要核验：点的映射由\trGlobalSections{}确定。$\Spec A$ 的一个点 $p$，可以等同于「在其上为零的整体函数所成的那个\trPrimeIdeal」。依练习~7.3.A，作为一个\trIdeal{}的 $\pi(p)$ 就是自然映射 $B\to\kappa(p)$ 的\trKernel。（这里我们用到了 $\pi$ 是\trLocallyRingedSpace{}之间的映射这一事实。）这正是「由环映射 $B\to A$ 诱导的集合映射 $\Spec A\to\Spec B$」当初被定义的方式（\S3.2.10）。

特别提请注意：若 $g\in B$，则 $\pi^{-1}(D(g))=D(\pi^{\sharp}g)$——这又一次用到了 $\pi$ 是\trLocallyRingedSpace{}的\trMorphism{}这一假设。

尚待证明的是：$\pi^{\sharp}\colon\OO_{\Spec B}\to\pi_*\OO_{\Spec A}$ 正是练习~7.2.F 所给出的那个\trSheaf{}的\trMorphism（对照练习~7.3.C）。在特异基底上核验它就足够了（\cref{ex:sheaf-morphisms-on-base}(a)）。现在我们想核验：对任何一个诱导出\trSheaf{}的映射 $\OO_{\Spec B}\to\pi_*\OO_{\Spec A}$ 的\trLocallyRingedSpace{}之间的映射，它在任一\trDistinguishedOpen{} $D(g)\subset\Spec B$（$g\in B$）上的截面映射，都由\trGlobalSections{}的映射 $B\to A$ 所确定。

考虑交换图
\[
\begin{tikzcd}[row sep=3.4em]
B \ar[r, equal] \ar[d, "\mathrm{res}_{\Spec B,\,D(g)}"] &
\Gamma(\Spec B,\OO_{\Spec B}) \ar[r, "\pi^{\sharp}{}_{\Spec B}"] &
\Gamma(\Spec A,\OO_{\Spec A}) \ar[r, equal] \ar[d, "\mathrm{res}_{\Spec A,\,D(\pi^{\sharp}g)}"'] &
A \\
B_g \ar[r, equal] &
\Gamma(D(g),\OO_{\Spec B}) \ar[r, "\pi^{\sharp}{}_{D(g)}"] &
\Gamma(D(\pi^{\sharp}g),\OO_{\Spec A}) \ar[r, equal] &
A_{\pi^{\sharp}g}=A\otimes_B B_g
\end{tikzcd}
\]
竖直的箭头（向诸\trDistinguishedOpen{}的限制）都是对 $g$ 的\trLocalization，于是下方的水平映射 $\pi^{\sharp}{}_{D(g)}$ 便由上方的映射所确定（它无非是对 $g$ 作\trLocalization）。\hfill$\square$

我们已经为下面的定义做好了准备。

\textbf{7.3.3.}\ \emph{\trhDefinition{}。}\quad 若 $X$ 与 $Y$ 是\trScheme，那么作为\trLocallyRingedSpace{}的\trMorphism{} $\pi\colon X\to Y$ 就称为一个\textbf{\trScheme{}的\trMorphism}（morphism of schemes）。至此我们便定义了\textbf{\trScheme{}的\trCategory}，记作 $\mathit{Sch}$。（于是我们也就有了\trIsomorphism{}的观念——与先前 \S4.3.1 处的含义相同——以及\trAutomorphism{}的观念。当我们把 $\pi$ 解读为「由 $Y$ 参数化的一族概形」时，$\pi$ 的\emph{\trTarget} $Y$ 有时被称为\textbf{\trBase\trScheme}（base scheme）或\textbf{\trBase}（base）；待我们在 \S10.3.2 中定义了\trMorphism{}的诸\trFiber{}之后，这话会更加清楚。）

用\trLocallyRingedSpace{}陈述的这条定义，很容易推出试行定义~7.2.2：

\textbf{7.3.D.（重要练习）}\quad 证明：\trScheme{}的\trMorphism{} $\pi\colon X\to Y$ 是一个「在局部看似诸\trAffineScheme{}的\trMorphism」的\trRingedSpace{}的\trMorphism。确切地说：若 $\Spec A$ 是 $X$ 的一个\trAffine\trOpenSubset，$\Spec B$ 是 $Y$ 的一个\trAffine\trOpenSubset，且 $\pi(\Spec A)\subset\Spec B$，则其所诱导的\trRingedSpace{}的\trMorphism{}是一个\trAffineScheme{}的\trMorphism。（为助你一臂之力，请注意：若 $W\subset X$ 与 $Y\subset Z$ 都是\trRingedSpace{}的\trOpenEmbedding，则任何\trRingedSpace{}的\trMorphism{} $X\to Y$ 都经复合 $W\to X\to Y\to Z$ 诱导出\trRingedSpace{}的一个\trMorphism{} $W\to Z$。）证明：只需在一组 $(\Spec A_i,\Spec B_i)$ 上核验即可，这里诸 $\Spec A_i$ 构成 $X$ 的一个\trOpenCover，诸 $\Spec B_i$ 构成 $Y$ 的一个\trOpenCover。

在实际操作中，我们将使用仿射覆盖式的解读，而把\trLocallyRingedSpace{}彻底抛诸脑后。特别地，说不严格一点：\trAffineScheme{}的\trCategory{}就是把箭头全部掉头的\trRing{}的\trCategory。更精确地说：

\textbf{7.3.E.（练习）}\quad 证明：\trRing{}的\trCategory{}与\trAffineScheme{}的\trOppositeCategory{}是等价的（关于\trOppositeCategory{}与\trEquivalenceOfCategories，可分别参看 \S1.1.16 与 \S1.1.21）。

特别地，从一个点到一个点的映射可以有许多个！比如，从点 $\Spec\mathbb{C}$ 到点 $\Spec\mathbb{C}$ 就有两个不同的映射：恒等映射（对应于恒等映射 $\mathbb{C}\to\mathbb{C}$），以及复共轭。（这样的映射甚至还有更多！）

显然（由关于\trLocallyRingedSpace{}的相应事实可知）：\trScheme{}的\trMorphism{}可粘（练习~7.3.B），且两个\trMorphism{}的\trComposition{}仍是\trMorphism。这个\trCategory{}中的\trIsomorphism{}，恰恰就是我们先前（\S4.3.1）定义它们时的样子。

\textbf{7.3.F.（富于启发的练习）}\quad （这道练习可以让你练一练「把概形切成\trAffine\trOpenSet{}来理解\trMorphism」的手艺。）请把下面这句话琢磨出道理来：「$\mathbb{A}^{n+1}_k\setminus\{\tilde{0}\}\to\mathbb{P}^n_k$ 由
\[
(x_0,x_1,\dots,x_n)\longmapsto[x_0,x_1,\dots,x_n]
\]
给出，是一个\trScheme{}的\trMorphism。」警示：你不能只说点往哪里去，还得说函数往哪里去。所以你也许得把 $\mathbb{A}^{n+1}_k\setminus\{\tilde{0}\}$ 切成一块块仿射的，逐一描述那些映射，并核验它们可粘。（把 $k$ 换成一般的\trRing{} $B$ 时，你还能推广吗？答案见\cref{ex:no-common-zeros-projective}。）

\textbf{7.3.4.\ 到\trAffineScheme{}的\trMorphism。}\quad 下面这项结果表明：要描述到一个\trAffineScheme{}的\trMorphism，无须费力气用诸\trAffine\trOpenSet{}把源头盖满，这很容易。

\textbf{7.3.G.（必作练习）}\quad 证明：诸\trMorphism{} $X\to\Spec A$ 与诸\trRingMorphism{} $A\to\Gamma(X,\OO_X)$ 天然地一一对应。提示：先证 $X$ 为\trAffine{}时如此；再用\trMorphism{}可粘这一事实（练习~7.3.B）。（这话在\trLocallyRingedSpace{}的\trCategory{}中甚至也成立。你尽可以在那样的一般性下证明它，但在\trScheme{}的\trCategory{}中要更容易。你或许会在这句更一般的陈述里，听到双星注记~6.2.5 那句格言的回声。）

特别地，从每个\trScheme{}到「其\trGlobalSections{}之\trRing{}的 $\Spec$」都存在一个\trCanonical{}的\trMorphism。（警告：即便 $X$ 是有限型的 $k$-\trScheme，\trGlobalSections{}之\trRing{}也可能很惹人厌！特别地，它未必\trFinitelyGenerated；见 \S19.11.13。）\trCanonical{}态射 $X\to\Spec\Gamma(X,\OO_X)$ 是\trIsomorphism，当且仅当 $X$ 是\trAffine{}的（即：\trIsomorphic{}于某个 $\Spec A$，其中 $A$ 是\trRing）；而且在此时，它正是注记~4.3.2 所暗示的那个\trIsomorphism。

\textbf{7.3.H.（容易的练习）}\quad 若 $S_\bullet$ 是一个\trFinitelyGenerated{}的分次 $A$-\trAlgebra，请描述一个自然的「\trStructureMorphism{}」（structure morphism）$\Proj S_\bullet\to\Spec A$。%

\textbf{7.3.5.}\ \emph{\trhRemark{}。}\quad 从必作练习~7.3.G 出发，只消一小步便能证明：概形的某些\trProduct{}是存在的：若 $A$ 与 $B$ 是\trRing，则 $\Spec A\times\Spec B=\Spec(A\otimes_{\mathbb{Z}}B)$；而若 $A$ 与 $B$ 是 $\mathbb{C}$-\trAlgebra，则 $\Spec A\times_{\Spec\mathbb{C}}\Spec B=\Spec(A\otimes_{\mathbb{C}}B)$。不过我们并不着急，且等到\mbox{练习~10.1.A}再正经讨论此事。

\textbf{7.3.6.}\ \extrahard\ \emph{供行家的旁注：$\Gamma$ 与 $\Spec$ 是\trAdjoint{}的。}\quad 我们有一个从\trRing{}到\trLocallyRingedSpace{}的\trContravariantFunctor{} $\Spec$，以及一个从\trLocallyRingedSpace{}到\trRing{}的\trContravariantFunctor{} $\Gamma$。事实上，$(\Gamma,\Spec)$ 竟是一个\trAdjointPair！（警示：我们只讨论过\trCovariantFunctor{}的伴随；倘若你当真在意，就得自己琢磨出如何为\trContravariantFunctor{}定义伴随。）如此一来，我们本可以通过径直要求 $\Spec$ 是 $\Gamma$ 的\trRightAdjoint{}来定义 $\Spec$。（一个有趣却不相干的旁问：如果你用的是\trRingedSpace{}而非\trLocallyRingedSpace，$\Gamma$ 依然有一个\trRightAdjoint。它是什么？）

我们能够轻易描述驶向\trAffineScheme{}的\trMorphism，这将使我们得以重访先前关于给定域或环上的概形的讨论（\S5.3.6）。

\textbf{7.3.7.\ \trComplex{}概形的范畴（或更一般地，$k$ 为\trField{}时诸 $k$-概形的范畴；或更一般地，$A$ 为\trRing{}时诸 $A$-概形的范畴；或更一般地，$S$ 为\trScheme{}时诸 $S$-概形的范畴）。}\quad $S$-\trScheme{}的\trCategory{} $\mathit{Sch}_S$（其中 $S$ 是一个\trScheme）定义如下。其\trObject（诸 $S$-概形）是形如
\[
\begin{tikzcd}[row sep=2.4em]
X \ar[d] \\
S
\end{tikzcd}
\]
的\trMorphism。（到 $S$ 的那个\trMorphism{}称为\textbf{\trStructureMorphism{}}。这一叫法的一个动机是：若 $S=\Spec A$，则\trStructureMorphism{}让 $X$ 的每个\trOpenSet{}上的函数都带有 $A$-\trAlgebra{}的结构；参看 \S5.3.6。）$S$-概形的\trCategory{}中的诸\trMorphism{}定义为交换图
\[
\begin{tikzcd}[row sep=2.4em, column sep=4em]
X \ar[r] \ar[d] & Y \ar[d] \\
S \ar[r, "="] & S
\end{tikzcd}
\]
它们更便利的写法是
\[
\begin{tikzcd}[row sep=2.4em, column sep=4em]
X \ar[r] \ar[dr] & Y \ar[d] \\
& S
\end{tikzcd}
\]
只要不引起混淆（即\trBase\trScheme{}已经自明），就只写出图中上面那一行。当 $S=\Spec A$（$A$ 为\trRing）时，我们便得到 $A$-\trScheme{}的观念；它与 \S5.3.6 中的那个定义是同一回事（练习~7.3.I），只是形式更令人满意。例如，复几何学家便可以考虑 $\mathbb{C}$-概形的范畴。

\textbf{7.3.I.（容易的练习）}\quad 证明：\S7.3.7 中给出的这个 $A$-\trScheme{}的定义，与 \S5.3.6 中较早的定义相一致。

\textbf{7.3.J.（容易的练习）}\quad 证明：$\Spec\mathbb{Z}$ 是\trScheme{}的\trCategory{}中的\trFinalObject。换言之，若 $X$ 是任何一个\trScheme，则存在唯一的到 $\Spec\mathbb{Z}$ 的\trMorphism。（因此，\trScheme{}的\trCategory{}同构于 $\mathbb{Z}$-\trScheme{}的\trCategory。）若 $A$ 是任何一个\trRing（例如某个\trField{} $k$），证明：$\Spec A$ 是 $A$-\trScheme{}的\trCategory{}中的\trFinalObject。

\textbf{7.3.8.\ \trScheme{}的\trCoproduct。}

\textbf{7.3.K.（容易的练习）}\quad 设 $X$ 与 $Y$ 是\trScheme。证明：二者的\trDisjointUnion{} $X\sqcup Y$（定义见\cref{ex:disjoint-union-schemes}）是\trScheme{}的\trCategory{}中 $X$ 与 $Y$ 的\trCoproduct，从而 \trCoproduct{}符号 $\sqcup$ 的使用便有了着落。（与\trIdempotent{}的一个出人意料的关联，见 \S3.6.3。）

\textbf{7.3.9.\ 从（某些）\trAffineScheme{}出发的\trMorphism。}\quad \emph{从}\trAffineScheme{}出发的\trMorphism，不像\emph{到}\trAffineScheme{}的\trMorphism{}那么简单，但有几种情形值得一提。

\vkexsync{L}
\begin{exercise}\label{ex:spec-local-ring-canonical-map}
(a)~设 $p$ 是\trScheme{} $X$ 的一个点。请描述一个\trCanonical{}（即无须选取的）\trMorphism{} $\Spec\OO_{X,p}\to X$。（提示：先对\trAffine{}的 $X$ 做这件事。但随后对一般的 $X$，务必证明你的\trMorphism{}与选取无关。）
(b)~定义一个\trCanonical{}的\trMorphism{} $\Spec\kappa(p)\to X$。（这个态射常被写作 $p\to X$；此时我们给 $p$ 一种显而易见的解读：把它看作\trScheme{} $\Spec\kappa(p)$。）
\end{exercise}

\textbf{7.3.M.（练习：从局部环的谱到 $X$ 的\trMorphism）}\quad 设 $X$ 是一个\trScheme，$(A,\mathfrak{m})$ 是一个\trLocalRing。设我们有一个\trScheme{}的\trMorphism{} $\pi\colon\Spec A\to X$，它把 $[\mathfrak{m}]$ 送到 $p$。证明：任何包含 $p$ 的\trOpenSet{}都包含 $\pi$ 的像。证明：$\Mor(\Spec A,X)$ 与
\[
\bigl\{(p\in X,\ \text{某个局部\trHomomorphism{}}\ \OO_{X,p}\to A)\bigr\}
\]
之间成双射。（可能的提示：\cref{ex:spec-local-ring-canonical-map}(a)。）

这些练习把我们引向\textbf{域\trValuedPoint{}}（field-valued points）的观念，进而是一般些的\textbf{环\trValuedPoint{}}（ring-valued points）的观念，再进而是一般些的\textbf{概形\trValuedPoint{}}（scheme-valued points）的观念。%

\textbf{7.3.10.\ \trhDefinition{}：\trFunctorOfPoints，与\trScheme{}的概形\trValuedPoint{}（以及环\trValuedPoint{}与域\trValuedPoint{}）。}\quad 若 $Z$ 是一个\trScheme，则\trScheme{} $X$ 的\textbf{$Z$-\trValuedPoint{}}（$Z$-valued points；记作 $X(Z)$）定义为诸映射 $Z\to X$。若 $A$ 是一个\trRing，则\trScheme{} $X$ 的\textbf{$A$-\trValuedPoint{}}（$A$-valued points；记作 $X(A)$）定义为该概形的诸 $(\Spec A)$-\trValuedPoint{}。（最常见的情形是 $A$ 为一个\trField。）

如果你在某个基底概形 $B$ 之上工作——例如，复代数几何学家将只考虑 $B=\Spec\mathbb{C}$ 上的概形与态射——那么在上面的定义里就暗含了一张结构映射 $Z\to B$（在 $X(A)$ 的情形则是 $\Spec A\to B$）。例如，对一位复几何学家而言，若 $X$ 是 $\mathbb{C}$ 上的概形，则 $X$ 的 $\mathbb{C}(t)$-\trValuedPoint{}对应于形如
\[
\begin{tikzcd}[row sep=2.6em, column sep=4em]
\Spec\mathbb{C}(t) \ar[r] \ar[dr, "\xi"'] & X \ar[d, "\pi"] \\
& \Spec\mathbb{C}
\end{tikzcd}
\]
的交换图，其中 $\pi\colon X\to\Spec\mathbb{C}$ 是 $X$ 的结构映射，而 $\xi$ 对应于环的显然包含 $\mathbb{C}\to\mathbb{C}(t)$。（警告：$k$-\trScheme{} $X$ 的一个 $k$-\trValuedPoint{}有时被称为 $X$ 的一个「\textbf{\trRationalPoint{}}」（rational point），这个叫法很危险，因为对世间大多数人而言，「有理」指的是 $\mathbb{Q}$。我们将使用更稳妥的说法：$X$ 的「$k$-\trValuedPoint{}」。另一个稳妥的选择是 $X$ 的「$k$-\trRationalPoint{}」或 $k$-点。）

「$Z$-\trValuedPoint」（以及 $A$-\trValuedPoint）这一叫法并不妥帖，因为我们先前已经定义过\trScheme{}的「点」这个观念，而 $Z$-\trValuedPoint（以及 $A$-\trValuedPoint）却（未必）是点！但这类用法在文献中早已根深蒂固。（不妨到索引的「点」（point）条目下去看看，修饰这个词的形容词还有更多彼此打架的用法。）

\textbf{7.3.11.}\quad 此外，我们将会看到「$Z$-\trValuedPoint{}的\trProduct」会如你所望的那样行事（\S10.1.3）。这套语言之所以能发人联想，还有一层相关的原因：记号 $X(Z)$ 暗示着把 $X$ 解读为一个从\trScheme{}到\trSet{}的\textbf{\trContravariantFunctor} $h_X$——即\trScheme{} $X$ 的（概形值）\trFunctorOfPoints（参看例~1.1.20）。

$A$-\trValuedPoint{}之所以有用，还有一个下里巴人式的理由：\emph{对诸 $f_i\in\mathbb{Z}[x_1,\dots,x_n]$（关系）而言，仿射概形 $\Spec\mathbb{Z}[x_1,\dots,x_n]/(f_1,\dots,f_r)$ 的 $A$-\trValuedPoint{}恰是方程组}
\[
f_1(x_1,\dots,x_n)=\cdots=f_r(x_1,\dots,x_n)=0
\]
\emph{在\trRing{} $A$ 中的解。}例如，$x^2+y^2=16$ 的有理解，恰是 $\Spec\mathbb{Z}[x,y]/(x^2+y^2-16)$ 的 $\mathbb{Q}$-\trValuedPoint{}；整数解恰是 $\mathbb{Z}$-\trValuedPoint{}。于是，仿射概形（在 $\mathbb{Z}$ 上\trFiniteType{}的）的 $A$-\trValuedPoint{}可以得到简单的解读。而在 $A$ 为局部环的特殊情形，一般概形的 $A$-\trValuedPoint{}也有很好的解读（\cref{ex:spec-local-ring-canonical-map} 与练习~7.3.M）。

\vkexsync{N}
\begin{exercise}\label{ex:valued-points-determine-morphisms}
(a)~（容易）证明：\trScheme{}的\trMorphism{} $X\to Y$ 诱导诸 $Z$-\trValuedPoint{}之间的一个映射 $X(Z)\to Y(Z)$。
(b)~注意：\trScheme{}的\trMorphism{} $X\to Y$ 并不由其「底下的」点的映射确定。（例子何在？）证明：它们\emph{确实}由其诱导的诸 $Z$-\trValuedPoint{}的映射所确定，其中 $Z$ 跑遍一切\trScheme。（提示：取 $Z=X$。在做这道练习的过程中，你将借着\mbox{练习~10.1.B}的外衣，把\trYonedaLemma{}大体证明出来。）
\end{exercise}

另一方面，\trProjectiveSpace{}的 $Z$-\trValuedPoint{}可能颇为微妙。有些映射我们能轻易写出——把下一道练习应用于 $X=\Spec A$（其中 $A$ 是一个 $B$-\trAlgebra）的情形，便可见一斑。

\begin{exercise}[容易（但出奇富于启发）的练习（参看练习~7.3.F）]\label{ex:no-common-zeros-projective}
(a)~设 $B$ 是一个\trRing。若 $X$ 是一个 $B$-\trScheme，而 $f_0,\dots,f_n$ 是 $X$ 上没有公共零点的 $n+1$ 个函数，证明：$[f_0,\dots,f_n]$ 给出一个 $B$-\trScheme{}的\trMorphism{} $X\to\mathbb{P}^n_B$。
(b)~设 $g$ 是 $X$ 上一个无处为零的函数，$f_i$ 如 (a) 中所述。证明：到 $\mathbb{P}^n_B$ 的诸\trMorphism{} $[f_0,\dots,f_n]$ 与 $[gf_0,\dots,gf_n]$ 是同一个。
\end{exercise}

\textbf{7.3.12.}\ \emph{\trhExample{}（重访练习~7.3.F）。}\quad 考虑 $\mathbb{A}^{n+1}$ 上的 $n+1$ 个函数 $x_0,\dots,x_n$（又名\trTrivialBundle{}（trivial bundle）的 $n+1$ 个\trSection）。它们在 $\mathbb{A}^{n+1}\setminus\{\tilde{0}\}$ 上没有公共零点。于是它们确定一个\trMorphism{} $\mathbb{A}^{n+1}\setminus\{\tilde{0}\}\to\mathbb{P}^n$。%

\textbf{7.3.13.}\quad 你或许会指望\cref{ex:no-common-zeros-projective}(a) 给出到射影空间（$B$ 上的）的一切\trMorphism。但事实并非如此。的确，就连恒等\trMorphism{} $X=\mathbb{P}^1_k\to\mathbb{P}^1_k$ 也不是这种形状，因为源头 $\mathbb{P}^1$ 手里没有非常值的\trGlobalSections{}可用来搭出这个映射。（源头为仿射的类似例子也有。）不过，\mbox{定理~15.2.2} 中有一个正确的推广（它刻画了从概形到射影概形的\emph{一切}映射）。这条结果粗略地讲：只要诸 $f_i$ 不是地道的函数，而是某个\trLineBundle{}的\trSection，这套想法就行得通。想要干干净净地理解到射影概形的映射，将成为理解\trLineBundle{}的一路重要动机。

下一节中我们将会见到更多描述到射影空间的映射的途径。\mbox{练习~15.2.A} 将给出另一种描述，它直接推广了\cref{ex:no-common-zeros-projective}(a)，并且将（在\mbox{定理~15.2.2} 中）被证明是一种「泛」的描述。

顺带一提：在\trGrothendieck{}之前，「该如何正确解读点带有环中『坐标』的射影空间」曾被当作一个实实在在的问题。这些困难源于函子式推理的缺席。而通往正确答案的线索其实早已存在（从\trManifold{}到 $\mathbb{RP}^n$ 的映射也会遭遇同样的问题）——如果你提出的是一个这样的几何问题（射影空间可是几何对象），答案就必然是几何的，而非纯代数的！

\textbf{7.3.14.\ \trMorphism{}的可视化：当\trNilpotent{}在场时为概形的映射绘图。}\quad 你现在已经知道如何让\trScheme{}的诸点（\S3.3）与\trNilpotent（\S4.2 与 \S6.6）在眼前显形。下面这道不精确的练习，将让你对「当\trNilpotent{}掺和进来时，如何为概形的映射绘图」有一点感觉。设 $a\in\mathbb{C}$。考虑由 $x\mapsto a\epsilon$ 给出的环映射 $\mathbb{C}[x]\to\mathbb{C}[\epsilon]/(\epsilon^2)$。回想：$\Spec\mathbb{C}[\epsilon]/(\epsilon^2)$ 可以画作一个带着\trTangentVector{}（tangent vector）的点（\S4.2）。若 $a\neq0$，你会如何为这个映射画像？若 $a=0$，你的图景又会如何改变？（此时\trTangentVector{}应当被「碾碎」。）

\mbox{练习~13.1.G} 将把这一点大幅推广。你不妨现在就去读读它的陈述，以期一乐。

\textbf{7.3.15.}\ \extrahard\ \textbf{供有相应背景的读者：\trComplex{}代数簇的\trAnalytification{}（analytification）。}\quad 警告：任何关于\trAnalytification{}的讨论，都只面向熟悉复解析簇这一观念的读者，以及有意愿沿着我们对概形的展开、自行平行地把它发展出来的读者。%

\textbf{7.3.P.（练习：把\trComplex{}代数簇\trAnalytification{}）}\quad 设 $X$ 是一个既约的有限型 $\mathbb{C}$-\trScheme。请定义相应的\textbf{复解析\trPrevariety{}}（complex analytic prevariety）$X_{\mathrm{an}}$。（解析\trPrevariety{}的定义，就是去掉\trHausdorff{}条件的\trVariety{}的定义。）警示：你的定义不得依赖于对 $X$ 的某个\trAffine\trCover{}的选取。（提示：先讲清楚如何把既约的有限型仿射 $\mathbb{C}$-概形\trAnalytification{}；然后粘合。）利用弱\trNullstellensatz（\cref{thm:hilbert-weak-nullstellensatz}），给出 $X$ 的诸\trClosedPoint{}与 $X_{\mathrm{an}}$ 的诸点之间的一个双射。（事实上，还可以构造出诸集合之间的一个\trContinuousMap{} $X_{\mathrm{an}}\to X$，它是\cref{ex:map-c2-a2q} 的推广。）%

\textbf{7.3.Q.（练习：\trAnalytification{}函子）}\quad 回想练习~7.3.P 中的\trAnalytification{}构造。对既约有限型 $\mathbb{C}$-概形的每个\trMorphism{} $\pi\colon X\to Y$（$\mathbb{C}$ 上的），定义复解析\trPrevariety{}之间的一个\trMorphism{} $\pi_{\mathrm{an}}\colon X_{\mathrm{an}}\to Y_{\mathrm{an}}$（即 $\pi$ 的\trAnalytification{}）。证明：\trAnalytification{}给出一个从既约有限型 $\mathbb{C}$-概形的\trCategory{}到复解析\trPrevariety{}的\trCategory{}的\trFunctor。

\textbf{7.3.16.}\ \emph{\trhRemark{}。}\quad 两个互不同构的\trVariety{}可以拥有同构的\trAnalytification{}。例如，\trSerre{}在复\trManifold{} $\mathbb{C}^*\times\mathbb{C}^*$ 上描述过两种不同的代数结构（见 [Ha2, p.~232] 与 [MO68421]\footnote{译注：[MO68421] 指 MathOverflow 问答第 68421 号。}）：一种是「显然的那一种」，另一种则是某条\trEllipticCurve{}上的一个 $\mathbb{P}^1$-丛再去掉一个\trSection。至于带有无穷多个代数结构的光滑复曲面，可参看 \S19.11.3。另一方面，紧致的复代数簇只能有一种代数结构（见 [Se3, \S19]）。

\textbf{7.3.17.}\ \emph{关于\trAnalytification{}的进一步事实。}\quad 欲知\trAnalytification{}的更多内容，可参看练习~11.2.K（\trHausdorff{}性）、\mbox{练习~14.1.J}、\mbox{练习~18.4.H}、\mbox{练习~18.4.U}（线丛）、\S18.6.5、\S20.1.9（次数）与\mbox{练习~21.2.W}（\trManifold）。复解析空间的背景，见 [GR]。
